\section{Movimientos de partículas en un fluido}
\begin{ejercicio} % // TODO: Ejercicio de Olga
    Tratamos de resolver el sistema:
    \begin{equation*}
        \left\{\begin{array}{l}
                \dot{x} = y \\
                \dot{y} = -x \\
                \dot{z} = 0
        \end{array}\right.
    \end{equation*}
    de donde obtenemos:
    \begin{align*}
        x(t) &= c_1\cos t + c_2\sen t \\
        y(t) &= -c_1\sen t +c_2\cos t \\
        z(t) &= c_3 \in \mathbb{R}
    \end{align*}
    para ciertas constantes $c_1,c_2,c_3\in \mathbb{R}$. El ejercicio consiste en demostrar que la curva $\alpha = (x,y,z)$ tiene como traza una circunferencia.\\

    \noindent
    Para ello, vemos por una parte que la curva $\alpha=(x,y,z)$ está contenida en el plano $P(z=3)$ y que esta misma cumple:
    \begin{equation*}
        x^2(t) + y^2(t) = c_1^2 + c_2^2 \qquad \forall t\in \mathbb{R}
    \end{equation*}
    Queda claro que $\tr \alpha$ queda contenida en la circunferencia de centro $0$ y radio $\sqrt{c_1^2+c_2^2}$. Veamos ahora que $\alpha$ recubre toda la circunferencia. Para ello, recordamos que todos los puntos de la circunferencia son de la forma:
    \begin{align*}
        x(t) &= r\cos(t+\theta) \\
        y(t) &= r\sen(t+\theta)
    \end{align*}
    para ciertos $r>0$ y $\theta\in \mathbb{R}$. Buscamos relacionar $r$ y $\theta$ con $c_1$ y $c_2$:
    \begin{align*}
        c_1\cos t + c_2 \sen t &= r\cos(t+\theta) \\
                               &= r(\cos t\cos\theta - \sen t\sen \theta) \\
                               &= r\cos t\cos\theta - r\sen t\sen\theta
    \end{align*}
    de donde:
    \begin{equation*}
        \left\{\begin{array}{l}
            c_1 = r\cos\theta \\
            c_2 = -r\sen\theta \\
            c_1^2 + c_2^2 = r^2(\cos^2\theta + \sen^2\theta) = r^2 \quad\Longrightarrow\quad r = \sqrt{c_1^2 + c_2^2}
        \end{array}\right. 
    \end{equation*}
    Despejando y suponiendo que $(c_1,c_2)\neq (0,0)$:
    \begin{equation}\label{eq:ej1}
       \cos\theta = \frac{c_1}{\sqrt{c_1^2+c_2^2}} , \qquad \sen\theta = \frac{-c_2}{\sqrt{c_1^2+c_2^2}}
    \end{equation}
    Vemos que el vector:
    \begin{equation*}
        \left(\frac{c_1}{\sqrt{c_1^2+c_2^2}}, \frac{-c_2}{\sqrt{c_1^2+c_2^2}}\right)
    \end{equation*}
    Tiene módulo 1, por lo que podemos encontrar $\theta\in \mathbb{R}$ verificando~\refeq{eq:ej1}.\\

    \noindent
    \textbf{Otra forma:} Buscamos ver que:
    \begin{equation*}
        \left\{\begin{array}{l}
                c_1\cos t + c_2\sen t = r\cos(t+\theta) \\
                -c_1\sen t + c_2\cos t = r\sen(t+\theta)
        \end{array}\right.
    \end{equation*}
    Tomando $r = \sqrt{c_1^2 + c_2^2}$ y $\theta\in \mathbb{R}$ verificando:
    \begin{equation*}
        \cos\theta = \frac{c_1}{\sqrt{c_1^2+c_2^2}}, \qquad \sen\theta = \frac{-c_2}{\sqrt{c_1^2+c_2^2}}
    \end{equation*}
    Vemos que se da la igualdad.
\end{ejercicio}

\begin{ejercicio}
    Sean $f_n:[0,1]\to \mathbb{R}$ dadas por:
    \begin{equation*}
        f_n(t) = \sqrt{t+\frac{1}{n}} \qquad \forall t\in [0,1], \quad \forall n\in \mathbb{N}
    \end{equation*}
    Probar que $\{f_n\}$ es equicontinua.\\

    \noindent
    Para ello, si miramos a la definición, buscamos acotar la diferencia:
    \begin{equation*}
        f_n(t) - f_n(s) = \sqrt{t+\frac{1}{n}} - \sqrt{s+\frac{1}{n}}
    \end{equation*}
    Por $\varepsilon$ sabiendo acotar $t-s$ por $\delta$. Si multiplicamos y dividimos por el conjugado:
    \begin{equation*}
        |f_n(t) - f_n(s)| = \frac{|t-s|}{\sqrt{t+\frac{1}{n}}+\sqrt{s+\frac{1}{n}}}
    \end{equation*}
    Ahora:
    \begin{equation*}
        |f_n(t) - f_n(s)| = \frac{|t-s|^{\nicefrac{1}{2}}}{\sqrt{t+\frac{1}{n}}+\sqrt{s+\frac{1}{n}}}|t-s|^{\nicefrac{1}{2}}
    \end{equation*}
    Tenemos que:
    \begin{equation*}
        |t-s| \leq t+s \quad\Longrightarrow\quad|t-s|^{\nicefrac{1}{2}} \leq \sqrt{t} + \sqrt{s}
    \end{equation*}
    De donde podemos acotar el primer término del producto anterior por 1, de donde:
    \begin{equation*}
        |f_n(t) - f_n(s)| \leq |t-s|^{\nicefrac{1}{2}} \qquad \forall t,s\in [0,1], \quad \forall n\geq 0
    \end{equation*}
    Ahora, para $\varepsilon>0$, tomando $\delta = \varepsilon^2$ tenemos que si $|t-s|<\delta$ entonces $|t-s|^{\nicefrac{1}{2}}<\varepsilon$, de donde:
    \begin{equation*}
        |f_n(t) - f_n(s)| \leq |t-s|^{\nicefrac{1}{2}} < \varepsilon \qquad \forall n\geq 0
    \end{equation*}
\end{ejercicio}

\begin{ejercicio}
    Probar que la sucesión anterior converge uniformemente.\\

    \noindent
    Tenemos que probar que:
    \begin{equation*}
        \forall \varepsilon>0~~\exists N\in \mathbb{N}~:~n\geq N \quad\Longrightarrow\quad |f_n(t) - f(t)| < \varepsilon \qquad \forall t\in [0,1]
    \end{equation*}
    Para cierta función $f:[0,1]\to \mathbb{R}$. En primer lugar, observamos que fijado $t\in [0,1]$ y usando que la raíz es continua obtenemos que:
    \begin{equation*}
        \{f_n(t)\} = \left\{\sqrt{t+\frac{1}{n}}\right\} \to \sqrt{t}
    \end{equation*}
    Por lo que si definimos $f:[0,1]\to \mathbb{R}$ como $f(t) = \sqrt{t}$ para todo $t\in [0,1]$ tenemos entonces que $\{f_n\}\to f$ puntualmente. Como la convergencia uniforme implica la puntual si la sucesión $\{f_n\}$ converge uniformemente, esta deberá hacerlo a $f$.\\

    \noindent
    Una vez tenemos claro cuál es la función $f$, tratamos de estudiar la siguiente diferencia con el fin de acotarla para $n\in \mathbb{N}$ suficientemente grande:
    \begin{align*}
        |f_n(t) - f(t)| &= \left|\sqrt{t+\frac{1}{n}}-\sqrt{t}\right| = \frac{\nicefrac{1}{n}}{\sqrt{t+\frac{1}{n}} + \sqrt{t}} \leq \frac{\nicefrac{1}{n}}{\nicefrac{1}{\sqrt{n}}} = \frac{1}{\sqrt{n}} \qquad \forall n\in \mathbb{N}
    \end{align*}
    Así, fijado $\varepsilon>0$, tenemos que existe $N\in \mathbb{N}$ de forma que $\nicefrac{1}{\sqrt{N}}< \varepsilon$, y si tomamos $n\in \mathbb{N}$ con $n\geq N$ obtendremos que:
    \begin{equation*}
        |f_n(t) - f(t)| \leq \frac{1}{\sqrt{n}} \leq \frac{1}{\sqrt{N}} < \varepsilon \qquad \forall t\in [0,1]
    \end{equation*}
    Acabamos de probar que $\{f_n\}\to f$ uniformemente.\\

    \noindent
    \textbf{Otra forma.} (ejercicio 2 de 2022) y tenía 3 apartados.
    \begin{enumerate}[label=\alph*)]
        \item Probar $\sqrt{a+b}\leq \sqrt{a}+\sqrt{b}\quad \forall a,b\geq 0$, como $x^2$ es creciente en $\mathbb{R}^2$ tenemos la desigualdad si y solo si:
            \begin{equation*}
                a+b \leq a+b+2\sqrt{ab} \quad\Longleftrightarrow\quad 0 \leq \sqrt{ab}
            \end{equation*}
        \item Hacer el ejercicio anterior, para el que podemos usar el apartado a).
            \begin{align*}
                |f_n(t) - f(t)| &= \sqrt{t+\frac{1}{n}} - \sqrt{t} \leq \sqrt{\frac{1}{n}}
            \end{align*}
        \item Se pregunta sobre la convergencia uniforme de la sucesión de las derivadas.

            Tenemos que:
            \begin{equation*}
                f_n'(t) = \frac{1}{2\sqrt{t+\frac{1}{n}}} \qquad \forall t\in [0,1]
            \end{equation*}
            Y para $t=0$ vemos que:
            \begin{equation*}
                f_n'(0) = \frac{\sqrt{n}}{2} \to \infty
            \end{equation*}
            Por lo que $\{f_n'(0)\}$ ni siquiera converge puntualmente, con lo que no puede converger uniformemente.
    \end{enumerate}
\end{ejercicio}

\begin{ejercicio} % // TODO: Ejercicio
    Probar que la única función que verifica:
    \begin{equation*}
        |f(t)-f(s)| \leq L|t-s|^\alpha
    \end{equation*}
    con $\alpha>1$ son las constantes.\\

    \noindent
    Sea $f:I\to \mathbb{R}$ una función verificando:
    \begin{equation*}
        |f(t)-f(s)| \leq L|t-s|^\alpha \qquad \forall t,s\in I
    \end{equation*}
    con $L>0$ y $\alpha>1$ tenemos entonces que para $h\neq 0$:
    \begin{equation*}
        \left|\frac{f(t+h)-f(t)}{h}\right| \leq L|h|^{\alpha-1}
    \end{equation*}
    Por lo que $\exists \lim\limits_{h\to0}\frac{f(t+h)-f(t)}{h} = 0$, luego $f$ tiene que ser constante.
\end{ejercicio}

\begin{ejercicio}
    ¿Existe una sucesión de funciones $\{f_n\}$ con $f_n:[0,1]\to \mathbb{R}$ continuas tales que $\{f_n\}$ converge uniformemente pero no es equicontinua? ¿Y si $f_n:\left]0,1\right[\to \mathbb{R}$?\\

    \noindent
    No:
    \begin{equation*}
        |f_n(t) - f_n(s)| \leq |f_n(t) - f(t)| + |f(t) - f(s)| + |f(s) - f_n(s)|
    \end{equation*}
    Las dos por convergencia uniforme y la de enmedio por continuidad uniforme, continua en compacto.\\

    \noindent
    Para el segundo sí que es cierto.
\end{ejercicio}

% // TODO: Ejercicio 5 del año 2025 1er parcial
\begin{ejercicio}
    Tenemos una sucesión de funciones $\{f_n\}_{n\geq 1}$ que vienen definidas como $f_n:[0,1]\to \mathbb{R}$ de modo que:
    \begin{equation*}
        |f_n(0)| \leq 7 \qquad \forall n\in \mathbb{N}
    \end{equation*}
    Con $\{f_n\}$ equicontinua y existe una parcial $\{f_{\sigma(n)}\}$ que converge uniformemente.\\

    \noindent
    Buscamos aplicar el Teorema de Arcoli, tenemos que ver que:
    \begin{equation*}
        |f_n(t)| \leq M \qquad \forall n\in \mathbb{N}, \quad \forall t\in [0,1]
    \end{equation*}
    Recordamos que $\{f_n\}$ es equicontinau si:
    \begin{equation*}
        \forall \varepsilon>0~~\exists \delta>0~:~ |f_n(t) - f_n(s)| < \varepsilon, \quad |t-s| < \delta\qquad \forall n\in \mathbb{N}
    \end{equation*}
    La equicontinuidad dice que la función ``crece por escaleras''. Vamos a probar la cota por inducción.\\

    \noindent
    Tomamos $\varepsilon = 1$ y la equicontinuidad nos da $\delta>0$, tomamos $m = \frac{\delta}{2}$, $k\geq 1$ y si $t\in [(k-1)m, km]\cap [0,1]\neq \emptyset $, vamos a probar que:
    \begin{equation*}
        |f_n(t)| \leq 7 + k \qquad \forall n\in \mathbb{N}
    \end{equation*}
    \begin{itemize}
        \item Para $k = 1$ tenemos $t\in [0,m]$, de donde:
            \begin{equation*}
                |f_n(t)| \leq |f_n(t)-f_n(0)| + |f_n(0)| \leq 1 + 7 = 8
            \end{equation*}
        \item Suponiendo ahora que $k > 1$, tomando $t\in [km, (k+1)m]$, tenemos que:
            \begin{equation*}
                |f_n(t)| \leq |f_n(t) - f_n(km)| + |f_n(km)| \leq 1 + 7 + k = 7 + (k+1)
            \end{equation*}
    \end{itemize}
    Como tenemos cota uniforme, aplicando el Teorema de Arcoli conseguimos el ejercicio.
\end{ejercicio}

% // TODO: Relacion 1 ej 2
\begin{ejercicio}
    % // TODO: Copiar el enunciado del ejercicio
    Tenemos $F\in C^1(\mathbb{R}^2)$ con la condición $F(x_0,y_0)=0$ y tenemos:
    \begin{equation*}
        \frac{\partial F}{\partial y}(x_0,y_0) \neq 0
    \end{equation*}
    % NO es problema de valores iniciales, es una ecuación funcional
    \begin{enumerate}[label=\alph*)]
        \item Tenemos que $\mathbb{R}^2$ es un abierto, que $F\in C^1(\mathbb{R}^2)$ y que:
            \begin{equation*}
                \frac{\partial F}{\partial y}(x_0,y_0) \neq 0
            \end{equation*}
            Aplicando el Teorema de la función implícita obtenemos directamente la existencia de la función.
        \item Queremos ahora llegar a la misma solución utilizando Cauchy-Peano. Para ello, hacemos derivación implícita, suponemos que tenemos una función $y$ que depende de $x$ y suponemos que tenemos una solución $F$, en cuyo caso se tendrá la ecuación:
            \begin{equation*}
                \frac{d}{dx}F(x,y) = \frac{\partial F}{\partial x}(x,y) + \frac{\partial F}{\partial y}(x,y)y' = 0
            \end{equation*}
            ¿Esta ecuación equivale a encontrar la $F$? tendríamos que $F(x,y(x)) = c$, de donde tendríamos la ecuación inicial que nos han dado. La idea es que dada una ecuación funcional se puede obtener una ecuación diferencial, pero en este caso podemos invertir este procedimiento. Aquí comienza la prueba, que va a consistir en aplicar el Teorema de Cauchy-Peano a la ecuación diferencial:

            \begin{equation*}
                y' = \frac{-\frac{\partial F(x,y)}{\partial y}}{\frac{\partial F}{\partial y}(x,y)}, \qquad y(x_0) = y_0
            \end{equation*}
            Como $F$ es de clase $C^1$ tenemos que es continua y como su parcial es distinta de cero, podemos exigir que existe un entorno $\Omega$ abierto y conexo de $(x_0,y_0)$ donde tengamos el denominador no nulo. Vemos que el campo es continuo como cociente de continuos, que $(x_0,y_0) \in \Omega^\circ$ y aplicando el Teorema de Cauchy-Peano conseguimos una solución $y:I\to \mathbb{R}$  de forma que $Gr(y)\subset \Omega$ y cumpliendo la ecuación diferencial. Ahora se continuaría hacia detrás.
    \end{enumerate}
\end{ejercicio}
